% GENERATED FILE. Edit canonical lessons, then run npm run generate:books.
% canonical-source-hash: fnv1a64:bb06208544c928d7
% canonical-lessons: HI-C443-nahar, HI-C443-jalasay, HI-C443-dhara, HI-C443-sicna, HI-C443-nirana

\chapter{Canals and Fields}
\label{ch:hi-canals-and-fields}
\hlchaptermodality{443}

\begin{chapteropening}
\textbf{By the end of this chapter:} \emph{I can talk about canals, reservoirs and currents, and watering and weeding the fields.}

Complete the last lesson of chapter 443: I can talk about canals, reservoirs and currents, and watering and weeding the fields.
\end{chapteropening}

\section[\texorpdfstring{\hi{नहर} (nahar)}{नहर (nahar)}]{\hi{नहर} (nahar) — a canal}

\label{lesson:HI-C443-nahar}



\begin{quote}
Before the new one: say the Hindi for a straw, then the Hindi for to roll (something).
\end{quote}

\subsection*{You'll want to know: \hi{नहर}}
\textbf{\hi{नहर}} — \emph{nahar} — \textquotedblleft{}a canal\textquotedblright{}.

\begin{quote}
\hi{इस} \hi{नहर} \hi{से} \hi{खेतों} \hi{को} \hi{पानी} \hi{मिलता} \hi{है।}
\end{quote}

\begin{quote}
\emph{is nahar se khetoṁ ko pānī miltā hai.}
\end{quote}

\begin{quote}
The fields get their water from this canal.
\end{quote}

Feminine.

\begin{cousinweb}
\textbf{\hi{नहर}} (\emph{nahar}) is usually traced to Arabic \emph{nahr}, \textquotedblleft{}a river, a stream\textquotedblright{}, through Persian; Hindi breaks the final hr with a vowel and uses it for a channel dug by people.
\end{cousinweb}

\begin{grammarlens}[title={What you've built}]
One more word for getting around.
\end{grammarlens}

\subsection*{Your turn}
\begin{itemize}
\raggedright
  \item \emph{Say it:} \emph{nahar}
  \item \emph{Say it:} \emph{nahar}, once more
  \item \emph{Say it:} \emph{is nahar se khetoṁ ko pānī miltā hai.} — The fields get their water from this canal.
  \item \emph{Recall:} say the Hindi for a straw, then the Hindi for to roll (something), then say \emph{nahar} again
\end{itemize}

\begin{tcolorbox}[breakable,colback=teal!4,colframe=teal!35!black,title={Before you move on}]
What does \hi{नहर} mean? (\textquotedblleft{}A canal\textquotedblright{}.) Say it once more.
\end{tcolorbox}
\section[\texorpdfstring{\hi{जलाशय} (jalāśay)}{जलाशय (jalāśay)}]{\hi{जलाशय} (jalāśay) — a reservoir}

\label{lesson:HI-C443-jalasay}



\begin{quote}
Before the new one: say the Hindi for to roll (something), then the Hindi for a canal.
\end{quote}

\subsection*{You'll want to know: \hi{जलाशय}}
\textbf{\hi{जलाशय}} — \emph{jalāśay} — \textquotedblleft{}a reservoir\textquotedblright{}.

\begin{quote}
\hi{गर्मी} \hi{में} \hi{जलाशय} \hi{का} \hi{पानी} \hi{बहुत} \hi{कम} \hi{हो} \hi{जाता} \hi{है।}
\end{quote}

\begin{quote}
\emph{garmī meṁ jalāśay kā pānī bahut kam ho jātā hai.}
\end{quote}

\begin{quote}
In the summer the water in the reservoir gets very low.
\end{quote}

Masculine.

\begin{cousinweb}
\textbf{\hi{जलाशय}} (\emph{jalāśay}) is a tatsama, Sanskrit \emph{jalāśaya}, \textquotedblleft{}a reservoir, a pond, a lake\textquotedblright{} (Monier-Williams), from \emph{jala}, \textquotedblleft{}water\textquotedblright{}, and \emph{āśaya}, \textquotedblleft{}a resting place\textquotedblright{}.
\end{cousinweb}

\begin{grammarlens}[title={What you've built}]
One more word for getting around.
\end{grammarlens}

\subsection*{Your turn}
\begin{itemize}
\raggedright
  \item \emph{Say it:} \emph{jalāśay}
  \item \emph{Say it:} \emph{jalāśay}, once more
  \item \emph{Say it:} \emph{garmī meṁ jalāśay kā pānī bahut kam ho jātā hai.} — In the summer the water in the reservoir gets very low.
  \item \emph{Recall:} say the Hindi for to roll (something), then the Hindi for a canal, then say \emph{jalāśay} again
\end{itemize}

\begin{tcolorbox}[breakable,colback=teal!4,colframe=teal!35!black,title={Before you move on}]
What does \hi{जलाशय} mean? (\textquotedblleft{}A reservoir\textquotedblright{}.) Say it once more.
\end{tcolorbox}
\section[\texorpdfstring{\hi{धारा} (dhārā)}{धारा (dhārā)}]{\hi{धारा} (dhārā) — a current; a stream}

\label{lesson:HI-C443-dhara}



\begin{quote}
Before the new one: say the Hindi for a canal, then the Hindi for a reservoir.
\end{quote}

\subsection*{You'll want to know: \hi{धारा}}
\textbf{\hi{धारा}} — \emph{dhārā} — \textquotedblleft{}a current; a stream\textquotedblright{}.

\begin{quote}
\hi{नदी} \hi{की} \hi{धारा} \hi{बहुत} \hi{तेज़} \hi{थी।}
\end{quote}

\begin{quote}
\emph{nadī kī dhārā bahut tez thī.}
\end{quote}

\begin{quote}
The river's current was very strong.
\end{quote}

Feminine.

\begin{cousinweb}
\textbf{\hi{धारा}} (\emph{dhārā}) is a tatsama, Sanskrit \emph{dhārā}, \textquotedblleft{}a stream or current of water\textquotedblright{} (Monier-Williams). In law and in the news it is also a section of an act.
\end{cousinweb}

\begin{grammarlens}[title={What you've built}]
One more word for getting around.
\end{grammarlens}

\subsection*{Your turn}
\begin{itemize}
\raggedright
  \item \emph{Say it:} \emph{dhārā}
  \item \emph{Say it:} \emph{dhārā}, once more
  \item \emph{Say it:} \emph{nadī kī dhārā bahut tez thī.} — The river's current was very strong.
  \item \emph{Recall:} say the Hindi for a canal, then the Hindi for a reservoir, then say \emph{dhārā} again
\end{itemize}

\begin{tcolorbox}[breakable,colback=teal!4,colframe=teal!35!black,title={Before you move on}]
What does \hi{धारा} mean? (\textquotedblleft{}A current; a stream\textquotedblright{}.) Say it once more.
\end{tcolorbox}
\section[\texorpdfstring{\hi{सींचना} (sī̃cnā)}{सींचना (sī̃cnā)}]{\hi{सींचना} (sī̃cnā) — to water; to irrigate}

\label{lesson:HI-C443-sicna}



\begin{quote}
Before the new one: say the Hindi for a reservoir, then the Hindi for a current.
\end{quote}

\subsection*{You'll want to know: \hi{सींचना}}
\textbf{\hi{सींचना}} — \emph{sī̃cnā} — \textquotedblleft{}to water; to irrigate\textquotedblright{}.

\begin{quote}
\hi{किसान} \hi{नहर} \hi{के} \hi{पानी} \hi{से} \hi{खेत} \hi{सींचते} \hi{हैं।}
\end{quote}

\begin{quote}
\emph{kisān nahar ke pānī se khet sī̃cte haiṁ.}
\end{quote}

\begin{quote}
The farmers water their fields from the canal.
\end{quote}

Past, with \hi{ने} (he or she …): \hi{उसने} \hi{बगीचा} \hi{सींचा}, \emph{usne bagīcā sī̃cā}.

\begin{cousinweb}
\textbf{\hi{सींचना}} (\emph{sī̃cnā}) is a tadbhava of Sanskrit \emph{siñcati}, \textquotedblleft{}pours out, sprinkles\textquotedblright{} (CDIAL 13394), which lists Hindi \emph{sī̃cnā}, \textquotedblleft{}to irrigate\textquotedblright{}; the nasal before c has become a nasal vowel.
\end{cousinweb}

\begin{grammarlens}[title={What you've built}]
One more word for getting around.
\end{grammarlens}

\subsection*{Your turn}
\begin{itemize}
\raggedright
  \item \emph{Say it:} \emph{sī̃cnā}
  \item \emph{Say it:} \emph{sī̃cnā}, once more
  \item \emph{Say it:} \emph{kisān nahar ke pānī se khet sī̃cte haiṁ.} — The farmers water their fields from the canal.
  \item \emph{Recall:} say the Hindi for a reservoir, then the Hindi for a current, then say \emph{sī̃cnā} again
\end{itemize}

\begin{tcolorbox}[breakable,colback=teal!4,colframe=teal!35!black,title={Before you move on}]
What does \hi{सींचना} mean? (\textquotedblleft{}To water; to irrigate\textquotedblright{}.) Say it once more.
\end{tcolorbox}
\section[\texorpdfstring{\hi{निराना} (nirānā)}{निराना (nirānā)}]{\hi{निराना} (nirānā) — to weed (a field)}

\label{lesson:HI-C443-nirana}



\begin{quote}
Before the new one: say the Hindi for a current, then the Hindi for to irrigate.
\end{quote}

\subsection*{You'll want to know: \hi{निराना}}
\textbf{\hi{निराना}} — \emph{nirānā} — \textquotedblleft{}to weed (a field)\textquotedblright{}.

\begin{quote}
\hi{किसान} \hi{धूप} \hi{में} \hi{खेत} \hi{निरा} \hi{रहे} \hi{थे।}
\end{quote}

\begin{quote}
\emph{kisān dhūp meṁ khet nirā rahe the.}
\end{quote}

\begin{quote}
The farmers were weeding the fields in the sun.
\end{quote}

Past, with \hi{ने} (he or she …): \hi{उसने} \hi{खेत} \hi{निराया}, \emph{usne khet nirāyā}.

\begin{cousinweb}
\textbf{\hi{निराना}} (\emph{nirānā}) goes back to \emph{*nīḍāti}, \textquotedblleft{}cuts down, weeds\textquotedblright{}, a form CDIAL reconstructs (CDIAL 7542), which lists Hindi \emph{nirānā}, \textquotedblleft{}to reap, clear weeds\textquotedblright{}.
\end{cousinweb}

\begin{grammarlens}[title={What you've built}]
One more word for getting around.
\end{grammarlens}

\subsection*{Your turn}
\begin{itemize}
\raggedright
  \item \emph{Say it:} \emph{nirānā}
  \item \emph{Say it:} \emph{nirānā}, once more
  \item \emph{Say it:} \emph{kisān dhūp meṁ khet nirā rahe the.} — The farmers were weeding the fields in the sun.
  \item \emph{Recall:} say the Hindi for a current, then the Hindi for to irrigate, then say \emph{nirānā} again
\end{itemize}

\begin{tcolorbox}[breakable,colback=teal!4,colframe=teal!35!black,title={Before you move on}]
What does \hi{निराना} mean? (\textquotedblleft{}To weed (a field)\textquotedblright{}.) Say it once more.
\end{tcolorbox}
